% GENERATED FILE. Edit canonical lessons, then run npm run generate:books.
% canonical-source-hash: fnv1a64:da59c0297eff9a9f
% canonical-lessons: ES-C529-explosion, ES-C529-explotar, ES-C529-estallar, ES-C529-escombros, ES-C529-destruir

\chapter{La explosión}
\label{ch:es-la-explosion}
\hlchaptermodality{529}

\begin{chapteropening}
\textbf{By the end of this chapter:} \emph{I can say that something exploded or burst, what it destroyed and what was left in the rubble.}

Complete the last lesson of chapter 529: I can say that something exploded or burst, what it destroyed and what was left in the rubble.
\end{chapteropening}

\section[la explosión]{la explosión — the explosion}

\label{lesson:ES-C529-explosion}



\begin{quote}
Before the new one: say the Spanish for the earthquake, then the Spanish for the flood.
\end{quote}

\subsection*{You'll want to know: la explosión}
\textbf{la explosión} — \textquotedblleft{}the explosion\textquotedblright{}.

\begin{quote}
Oímos una explosión muy fuerte.
\end{quote}

\begin{quote}
\emph{We heard a very loud explosion.}
\end{quote}

\begin{cousinweb}
\emph{explosión} is a learned borrowing from Latin \emph{explōsiō} 'a driving off the stage by clapping' (Lewis \& Short), built on \emph{plaudere} 'to clap'. The noise of the crowd became the noise of a blast.
\end{cousinweb}

\begin{grammarlens}[title={What you've built}]
One more word for telling, explaining and describing at B1.
\end{grammarlens}

\subsection*{Your turn}
\begin{itemize}
\raggedright
  \item \emph{Say it:} \emph{la explosión}
  \item \emph{Say it:} \emph{la explosión}, once more
  \item \emph{Say it:} a sentence of your own with \emph{la explosión}
  \item \emph{Recall:} say the Spanish for the earthquake, then the Spanish for the flood, then say \emph{la explosión} again
\end{itemize}

\begin{tcolorbox}[breakable,colback=teal!4,colframe=teal!35!black,title={Before you move on}]
What does la explosión mean? (\textquotedblleft{}The explosion\textquotedblright{}.) Say it once more.
\end{tcolorbox}
\section[explotar]{explotar — to explode; to exploit}

\label{lesson:ES-C529-explotar}



\begin{quote}
Before the new one: say the Spanish for the flood, then the Spanish for the explosion.
\end{quote}

\subsection*{You'll want to know: explotar}
\textbf{explotar} — \textquotedblleft{}to explode; to exploit\textquotedblright{}.

\begin{quote}
La bombona de gas explotó.
\end{quote}

\begin{quote}
\emph{The gas canister exploded.}
\end{quote}

It also means to exploit: \emph{explotar a los trabajadores}.

\begin{cousinweb}
\emph{explotar} was borrowed from French \emph{exploiter} 'to make use of' (EtymDB). The sense 'to explode' is usually explained as a later Spanish growth, pulled along by \emph{explosión}.
\end{cousinweb}

\begin{grammarlens}[title={What you've built}]
One more word for telling, explaining and describing at B1.
\end{grammarlens}

\subsection*{Your turn}
\begin{itemize}
\raggedright
  \item \emph{Say it:} \emph{explotar}
  \item \emph{Say it:} \emph{explotar}, once more
  \item \emph{Say it:} a sentence of your own with \emph{explotar}
  \item \emph{Recall:} say the Spanish for the flood, then the Spanish for the explosion, then say \emph{explotar} again
\end{itemize}

\begin{tcolorbox}[breakable,colback=teal!4,colframe=teal!35!black,title={Before you move on}]
What does explotar mean? (\textquotedblleft{}To explode; to exploit\textquotedblright{}.) Say it once more.
\end{tcolorbox}
\section[estallar]{estallar — to burst, to break out}

\label{lesson:ES-C529-estallar}



\begin{quote}
Before the new one: say the Spanish for the explosion, then the Spanish for to explode.
\end{quote}

\subsection*{You'll want to know: estallar}
\textbf{estallar} — \textquotedblleft{}to burst, to break out\textquotedblright{}.

\begin{quote}
Estalló una tormenta por la tarde.
\end{quote}

\begin{quote}
\emph{A storm broke out in the afternoon.}
\end{quote}

Also for wars and scandals: \emph{Estalló la guerra.}

\begin{cousinweb}
\emph{estallar} is usually traced to an old verb \emph{*astellar} 'to break into splinters', built on \emph{astilla} 'a splinter'; the details are not settled.
\end{cousinweb}

\begin{grammarlens}[title={What you've built}]
One more word for telling, explaining and describing at B1.
\end{grammarlens}

\subsection*{Your turn}
\begin{itemize}
\raggedright
  \item \emph{Say it:} \emph{estallar}
  \item \emph{Say it:} \emph{estallar}, once more
  \item \emph{Say it:} a sentence of your own with \emph{estallar}
  \item \emph{Recall:} say the Spanish for the explosion, then the Spanish for to explode, then say \emph{estallar} again
\end{itemize}

\begin{tcolorbox}[breakable,colback=teal!4,colframe=teal!35!black,title={Before you move on}]
What does estallar mean? (\textquotedblleft{}To burst, to break out\textquotedblright{}.) Say it once more.
\end{tcolorbox}
\section[los escombros]{los escombros — the rubble, the debris}

\label{lesson:ES-C529-escombros}



\begin{quote}
Before the new one: say the Spanish for to explode, then the Spanish for to burst.
\end{quote}

\subsection*{You'll want to know: los escombros}
\textbf{los escombros} — \textquotedblleft{}the rubble, the debris\textquotedblright{}.

\begin{quote}
Sacaron a dos personas de entre los escombros.
\end{quote}

\begin{quote}
\emph{They pulled two people out of the rubble.}
\end{quote}

Almost always in the plural.

\begin{cousinweb}
\emph{escombros} is built on \emph{escombrar} 'to clear out', which EtymDB traces to spoken Latin \emph{*excomborāre} and a Celtic word \emph{*comboros}, usually glossed 'a barrier, a heap'.
\end{cousinweb}

\begin{grammarlens}[title={What you've built}]
One more word for telling, explaining and describing at B1.
\end{grammarlens}

\subsection*{Your turn}
\begin{itemize}
\raggedright
  \item \emph{Say it:} \emph{los escombros}
  \item \emph{Say it:} \emph{los escombros}, once more
  \item \emph{Say it:} a sentence of your own with \emph{los escombros}
  \item \emph{Recall:} say the Spanish for to explode, then the Spanish for to burst, then say \emph{los escombros} again
\end{itemize}

\begin{tcolorbox}[breakable,colback=teal!4,colframe=teal!35!black,title={Before you move on}]
What does los escombros mean? (\textquotedblleft{}The rubble, the debris\textquotedblright{}.) Say it once more.
\end{tcolorbox}
\section[destruir]{destruir — to destroy}

\label{lesson:ES-C529-destruir}



\begin{quote}
Before the new one: say the Spanish for to burst, then the Spanish for the rubble.
\end{quote}

\subsection*{You'll want to know: destruir}
\textbf{destruir} — \textquotedblleft{}to destroy\textquotedblright{}.

\begin{quote}
El fuego destruyó tres casas.
\end{quote}

\begin{quote}
\emph{The fire destroyed three houses.}
\end{quote}

It takes a \emph{y}: \emph{destruyo}, \emph{destruyó}.

\begin{cousinweb}
\emph{destruir} comes from Latin \emph{dēstruere} 'to tear down, to destroy' (EtymDB), the reverse of \emph{struere} 'to pile up, to build'.
\end{cousinweb}

\begin{grammarlens}[title={What you've built}]
One more word for telling, explaining and describing at B1.
\end{grammarlens}

\subsection*{Your turn}
\begin{itemize}
\raggedright
  \item \emph{Say it:} \emph{destruir}
  \item \emph{Say it:} \emph{destruir}, once more
  \item \emph{Say it:} a sentence of your own with \emph{destruir}
  \item \emph{Recall:} say the Spanish for to burst, then the Spanish for the rubble, then say \emph{destruir} again
\end{itemize}

\begin{tcolorbox}[breakable,colback=teal!4,colframe=teal!35!black,title={Before you move on}]
What does destruir mean? (\textquotedblleft{}To destroy\textquotedblright{}.) Say it once more.
\end{tcolorbox}
